% §2 What is used from the earlier papers --- the recap of the two released modules of Paper V and
% of the isotropic paper this paper rests on, and the cited facts it spends (CHANGED Q-0281 wave 2: blind review W2F-2).
% Five since round 3 (C5).
% Written from the
% draft's §2 (Q-0285, 2026-10-04), on Paper V's model (paper-b/02-line-paper.tex,
% paper-c/02-earlier-modules.tex): the paper is read on its own, so what the later sections read is
% RESTATED here, cited by released version.
%
% Q-0307 (the split at §3 | §4, adr/0001). The single article's section 2 served its section 3 as
% well; that part went with the isotropic paper (paper-iso/02-line-papers.tex): the causal bridge,
% the dimension filtration, the extreme rays, the origin on the line, the two facts of the line used
% in proofs, and the cited facts A1, A2 and A4. Kept here: the notation, the axioms and the
% classification on the line (2.1, 2.2), and A3, A5 and A6 (2.4), verbatim. New: 2.3, the four (six since Q-0317, with the two marginal lemmas)
% statements of the isotropic paper that sections 3-4 and the appendices use, RESTATED as Paper V's
% modules restate the line paper (paper-b/02-line-paper.tex): verbatim from the blueprint under the
% same `% shared with blueprint` marker, so that `linkage check` guards them too, with [Iso]'s
% statement number in the title and the proof left there (`% proof omitted`). The statement lives in
% the isotropic paper, which proves it; here it is a restatement (article-kit LINKAGE.md, "A label,
% once per module"). Q-0457: [Iso] is released (iso-v0.1, 2026-10-14, version DOI
% 10.5281/zenodo.23357468, concept DOI 10.5281/zenodo.23357467) and cited through the librarian's
% citekey fagerstrom2026isotropic by its printed numbers, read from the released PDF (the export
% github.com/danielfagerstrom/spatial-hemigroup-scale-space-isotropic, tag v0.1): Definition 3.2,
% Lemmas 3.4 and 3.5, Theorem 3.6, Definition 3.10 and Lemma A.2, which are all in the paper and none
% in its supplement. Until then it was cited by the name of the statement (Q-0312, review flag 2),
% with a release marker at each pointer (records/affine/NODES-q0457.md lists them).
%
% No blueprint node in 2.2 and 2.4: this repository's blueprint has no node for a statement of [V]
% or [B] (repo CLAUDE.md: no \uses edge crosses into another repository's blueprint), so every such
% statement below carries a `% no blueprint node:` opt-out naming what it restates (article-kit
% LINKAGE.md rule 12). Each restatement was checked against the released source at its tag (v0.1:
% paper/02, 04, 06, 07), not against the draft. Theorem 2.3 is [V, Thm. 7.3] as released (with its
% uniqueness clause); the draft's last sentence ("the admissible exponents are exactly the exponents
% of the symmetric self-decomposable laws") is [V, Prop. 2.13], and is prose after the theorem. The
% cited facts are shared with their blueprint interface nodes thm:cited-A<N> (the author's decision
% of 2026-10-06), verbatim or pinned; those nodes are the ledger's "statement as used"
% (blueprint/AXIOMS.md), which was read at the page, and the page anchors are the ledger's. Module C is not cited (the author, 2026-10-04).

% Round 2 (presentation review), item 1: the sentence that re-described the three earlier papers is
% cut, since section 1's paragraph "The axioms" introduces them
% (records/affine/reviews/VERIFY-round2-front-paper.md, item 1, the safe part).
\section{What is used from the earlier papers}\label{sec:line}

This section restates, without proof, the statements of the line paper
\sscite{fagerstrom2026spatial} and of the isotropic paper \sscite{fagerstrom2026isotropic} that the
later sections use in their statements or in more than one proof, with the notions of the cone
paper \sscite{fagerstrom2026cone} that they read, and states the five theorems of the literature
that the proofs cite. The three papers are cited in their released
version~0.1, by the statement numbers of that version. A restated statement carries its released
number in its title, and its proof is at that number; a statement used once, inside one proof, is
restated in that proof. With the restatements the statements of this paper can be read without
the earlier papers, and its proofs rest on theirs. A reader who knows the earlier papers can go to
\S\ref{sec:cited-facts}.

\subsection{Notation}\label{sec:notation}

% Round 2 (presentation review), item 1: the vocabulary sentences, which repeated section 1's
% paragraph "The axioms", are cut to one sentence and the formula. Item 15's "or as a stage" belongs
% to the sentence on "kernel" that remains in section 1 only.
As in \S\ref{sec:introduction}, a probability measure is a \emph{law}, and its transform is its
characteristic function, $\widehat\mu(\xi) = \int e^{i\langle\xi,x\rangle}\,\mu(dx)$ for a law $\mu$
on $\RR^d$. For $A \in \GL(d)$, $A_*\mu := \mu \circ A^{-1}$ is the image measure, with
$\widehat{A_*\mu}(\xi) = \widehat\mu(A^{\tp}\xi)$. On functions $A$ acts by
$D_Af := |\det A|^{-1}f(A^{-1}\,\cdot)$, which preserves the integral, so that
$D_A(\mu * f) = (A_*\mu) * D_Af$. On the line the dilation by $\lambda > 0$ is
$D_\lambda f = \lambda^{-1}f(\cdot/\lambda)$, the translation by $a$ is $\tau_af = f(\cdot - a)$ and
the reflection is $Rf = f(-\,\cdot)$. $\Sym(d)$ is the space of real symmetric $d \times d$
matrices, $\PSD(d)$ and $\PD(d)$ the positive semidefinite and the positive definite ones, and
$\Sigma \le \Sigma'$ the Loewner order, $\Sigma' - \Sigma \in \PSD(d)$. $\Orth(d)$ is the
orthogonal group, $S^{d-1}$ the unit sphere, $\sigma_{d-1}$ its surface measure and
$|S^{d-1}| = \sigma_{d-1}(S^{d-1})$ its surface area. The frequency has three letters, each
following its source. From \S\ref{sec:classical} on, $\omega$ is the frequency on $\RR^d$, and a
proof that needs a second frequency writes it $\xi$. The
statements of the isotropic paper restated in \S\ref{sec:iso-recap} write it $\xi$, as that paper does,
and the frequency on the line $\omega$. The cited fact C2 writes it $z$, as its source does.
Table~\ref{tab:notation}, at the end of this section, collects the symbols of
\S\S\ref{sec:classical}--\ref{sec:landscape}, each with the statement that defines it, and the list
after it names the letters that keep a second meaning.

\subsection{The axioms and the classification on the line}\label{sec:line-axioms}

The line paper's object is a family of operators on $L^1(\RR)$ indexed by pairs of scales
$s \le t$. The operator $\Phi_{s,t}$ takes the record of a signal at scale $s$ to its record at the
coarser scale $t$. The axioms are these.

% no blueprint node: restates Definition 3.1 of [V] (fagerstrom2026spatial, v0.1)
\begin{definition}[symmetric cascade measurement family, {\prog{V}{Def.~3.1}}]\label{def:line-family}
A family $(\Phi_{s,t})_{0 \le s \le t}$ of linear operators on $L^1(\RR)$ is a \emph{symmetric
cascade measurement family} if:
\begin{enumerate}
\item[(A1)] each $\Phi_{s,t}$ is a bounded operator on $L^1(\RR)$;
\item[(A2)] $\Phi_{s,t}\tau_a = \tau_a\Phi_{s,t}$ for every $a \in \RR$;
\item[(A3)] $\Phi_{s,t}R = R\,\Phi_{s,t}$;
\item[(A4)] $\Phi_{s,t}$ maps nonnegative functions to nonnegative functions;
\item[(A5)] $\int\Phi_{s,t}f = \int f$ for every nonnegative $f \in L^1(\RR)$;
\item[(A6)] $\Phi_{t,t} = \mathrm{Id}$, and $\Phi_{s,t}\Phi_{r,s} = \Phi_{r,t}$ whenever
$r \le s \le t$;
\item[(A7)] for every $f \in L^1(\RR)$ the map $(s,t) \mapsto \Phi_{s,t}f$ is continuous from
$\{0 \le s \le t\}$ to $L^1(\RR)$;
\item[(A8)] there is a family $(S_\lambda)_{\lambda > 0}$ of increasing bijections of
$[0,\infty)$ with $D_\lambda\Phi_{s,t} = \Phi_{S_\lambda s,\,S_\lambda t}D_\lambda$ for all
$0 \le s \le t$ and all $\lambda > 0$;
\item[(ND)] $\Phi_{s,t} \ne \mathrm{Id}$ whenever $s < t$.
\end{enumerate}
\end{definition}

Axiom (A6), the \emph{cascade}, requires
that the stages compose, and leaves free how the stage from $s$ to $t$ depends on the two scales:
the family need not be a semigroup in any scale parameter. Axiom (A8) is the scale invariance of
the theory. Dilating the signal changes which scale a stage belongs to, and (A8) requires only that it
relabel the scales by an increasing bijection $S_\lambda$, without fixing in advance how. Axiom (ND)
excludes stages that do nothing. The first consequence of the axioms is that every stage is a
convolution.

% no blueprint node: restates Lemma 4.1 of [V] (fagerstrom2026spatial, v0.1)
\begin{lemma}[from operators to kernels on the line, {\prog{V}{Lem.~4.1}}]\label{lem:line-representation}
Under (A1), (A2), (A4) and (A5), for each pair $0 \le s \le t$ there is a unique probability
measure $\nu_{s,t}$ on $\RR$ with $\Phi_{s,t}f = \nu_{s,t} * f$ for every $f \in L^1(\RR)$; under
(A3) it is symmetric. Conversely every symmetric probability measure on $\RR$ defines an operator
satisfying (A1)--(A5).
\end{lemma}

The same passage in dimension $d$, for every group and order of the scheme of
\S\ref{sec:classical}, is Lemma~\ref{lem:operators-to-kernels}, and its converse is
Lemma~\ref{lem:kernels-to-operators}; both are proved in Appendix~\ref{app:classical}.

A \emph{gauge} is an increasing bijection $\chi$ of $[0,\infty)$, hence continuous, with
$\chi(0) = 0$. It labels the scales: the classification below describes a family by a gauge
and an exponent, and in a family with gauge $\chi$ axiom (A8) holds with
$S_\lambda = \chi^{-1}(\lambda\,\chi(\cdot))$, so that a dilation by $\lambda$ multiplies $\chi$ by
$\lambda$. A function
$F : [0,\infty) \to [0,\infty)$ is \emph{admissible} if it has the form \prog{V}{(2.3)}
\begin{equation}\label{eq:admissible-line}
  F(r) = a_{\mathrm V}\,r^2 + \int_0^\infty \bigl(1 - \cos rx\bigr)\,k(x)\,\frac{dx}{x},
\end{equation}
with $a_{\mathrm V} \ge 0$ and $k : (0,\infty) \to [0,\infty)$ nonincreasing,
$\int_0^1 x\,k(x)\,dx < \infty$ and $\int_1^\infty k(x)\,x^{-1}dx < \infty$. The number
$a_{\mathrm V}$ is the \emph{Gaussian coefficient} and $k$ the \emph{folded profile}. A subscript
$\mathrm V$ marks an object of the line, and an upright subscript $\mathrm I$, in
\S\ref{sec:iso-recap}, one of the time-causal theory; the italic $I$ of $\mu_I$ is the identity
matrix. A law is
\emph{infinitely divisible} when, for every $n$, it is the $n$-fold convolution of some law $\mu_n$
with itself. An admissible $F$ is the exponent of a symmetric infinitely divisible law, and the
measure $k(x)\,dx/x$ on $(0,\infty)$ is that law's \emph{\Levy{} measure} folded onto the half-line.
Displacements of size about $x$ arrive at rate $k(x)\,dx/x$, the two directions counted together. The two integrability conditions are
those of \prog{V}{Lem.~2.15}.

% no blueprint node: restates Theorem 7.3 of [V] (fagerstrom2026spatial, v0.1)
\begin{theorem}[the classification on the line, {\prog{V}{Thm.~7.3}}]\label{thm:line-classification}
A family $(\Phi_{s,t})_{0 \le s \le t}$ of operators on $L^1(\RR)$ satisfies (A1)--(A8) and (ND)
if and only if there are a gauge $\chi$ and an admissible $F \not\equiv 0$ such that
$\Phi_{s,t}f = \nu_{s,t} * f$ with symmetric probability measures $\nu_{s,t}$ on $\RR$,
\[
  \widehat\nu_{s,t}(\omega) = \exp\bigl[-\bigl(F(\chi(t)\,\omega) - F(\chi(s)\,\omega)\bigr)\bigr] .
\]
The pair $(\chi, F)$ is unique up to the normalization $\chi(1) = 1$.
\end{theorem}

The uniqueness says that $(c\,\chi, F(c^{-1}\,\cdot))$ gives the same family for every $c > 0$, and
nothing else does; the normalization $\chi(1) = 1$ fixes $c$. In the \emph{canonical gauge}, $\chi$
the identity, the kernel from scale zero at scale $t$ has
transform $e^{-F(t\omega)}$: it is the dilation by $t$ of one law. A law
$\mu$ on $\RR^d$ is \emph{self-decomposable} if for every $c \in (0,1)$ there is a law $\rho_c$, the
\emph{residual}, with
\begin{equation}\label{eq:self-decomposable}
  \widehat\mu(\xi) = \widehat\mu(c\,\xi)\,\widehat\rho_c(\xi), \qquad \xi \in \RR^d
\end{equation}
(\sscite{sato1999levy}, Def.~15.1, p.~90; on the line \prog{V}{Def.~2.12}). In words, the law is a
contracted copy of itself convolved with an independent remainder, for every contraction ratio.
This is what makes a cascade along one scale direction possible: the residuals are the stages.
The admissible functions are exactly the exponents of the symmetric self-decomposable laws on the
line \prog{V}{Prop.~2.13}.

\subsection{What is used from the isotropic paper}\label{sec:iso-recap}

The isotropic paper measures an image on $\RR^d$ with apertures that have a size and no shape and
no preferred direction. In this paper its theory is the isotropic corner of the scheme
(Proposition~\ref{prop:isotropic-corner}), and the isotropic member $\mu_I$ of an affine orbit is
described in its terms (\S\ref{sec:affine-stable-cone}). The later sections use six of its
statements: the axioms, the classification with the two lemmas that reduce it to the line, the
description of a law by its polar profile, and one estimate of the exponent. The isotropic paper
states its axioms on operators, as the line's with $\RR^d$ in place of $\RR$ and the orthogonal
group in place of the reflection \prog{Iso}{Def.~3.1}, and on kernels, the form restated here; in
this paper's scheme Lemmas~\ref{lem:operators-to-kernels} and~\ref{lem:kernels-to-operators} join
the two.

% shared with blueprint def:isotropic-kernel-family
\begin{definition}[isotropic cascade family, {\prog{Iso}{Def.~3.2}}]\label{def:isotropic-kernel-family}
Fix an integer $d \ge 1$. For $A \in \GL(d)$ write $A_*\mu := \mu \circ A^{-1}$ for the image of a
measure $\mu$ on $\RR^d$, and $\widehat\mu(\xi) = \int e^{i\langle\xi,x\rangle}\mu(dx)$, so that
$\widehat{A_*\mu}(\xi) = \widehat\mu(A^{\tp}\xi)$. An \emph{isotropic cascade family in dimension
$d$} is a family $(\mu_{s,t})_{0 \le s \le t}$ of probability measures on $\RR^d$ such that
\begin{enumerate}
\item[(I1)] \emph{(cascade)} $\mu_{t,t} = \delta_0$, and $\mu_{r,t} = \mu_{s,t} * \mu_{r,s}$
whenever $r \le s \le t$;
\item[(I2)] \emph{(continuity)} for every $\xi \in \RR^d$ the map
$(s,t) \mapsto \widehat\mu_{s,t}(\xi)$ is continuous on $\{0 \le s \le t\}$;
\item[(I3)] \emph{(rotation invariance)} $Q_*\mu_{s,t} = \mu_{s,t}$ for every $Q \in \Orth(d)$
and every pair;
\item[(I4)] \emph{(dilation covariance)} there is a family $(S_\lambda)_{\lambda > 0}$ of
increasing bijections of $[0,\infty)$ with $(\lambda I)_*\mu_{s,t} = \mu_{S_\lambda s,\,S_\lambda t}$
for every $\lambda > 0$ and every pair;
\item[(ND)] \emph{(nondegeneracy)} $\mu_{s,t} \ne \delta_0$ whenever $s < t$.
\end{enumerate}
\end{definition}

The cone paper names the class $\mathcal{A}_d$ of the admissible $F$ of \ref{eq:admissible-line}
for which $\xi \mapsto F(|\xi|)$, $\xi \in \RR^d$, is the exponent of a self-decomposable law on
$\RR^d$ \prog{B}{Prop.~4.9}. By the next theorem its members other than $0$ are the exponents of
the isotropic families in $\RR^d$. The Gaussian exponent $a_{\mathrm V}r^2$ lies in every
$\mathcal A_d$.

% shared with blueprint thm:isotropic-classification
\begin{theorem}[the isotropic classification in dimension $d$, {\prog{Iso}{Thm.~3.6}}]\label{thm:isotropic-classification}
% Q-0281 wave 1 (group E): proved. Round 2 (presentation review), item 12: this comment is brought
% up to date (it described the include line that Q-0365 removed). Since Q-0365 (2026-10-08)
% AffineHemigroup.isotropic_classification prints, besides Lean core, this module's ledger
% A3(ii)-(iii) (AffineHemigroup.levy_khintchine_unique and levy_khintchine_converse), through the
% necessity direction: ScaleSpaceCore's ScaleSpace.main_analysis with its hypothesis SymLevyUnique
% discharged by AffineHemigroup.line_levy_unique (prop:line-levy-unique). It prints no interface of
% Paper V's, and blueprint/trust-boundary.txt has no include line. The node is [A] on ledger A3. The
% declaration is guarded in CIAxiomGuard.lean; the other declarations of the tag, and the
% sufficiency direction, are Lean core and guarded.
Let $d \ge 1$. A family $(\mu_{s,t})_{0 \le s \le t}$ of probability measures on $\RR^d$ is an
isotropic cascade family in dimension $d$ if and only if there are a gauge $\chi$ and an
$F \in \mathcal{A}_d$ with $F \not\equiv 0$ such that
\begin{equation}\label{eq:isotropic-classification}
  \widehat\mu_{s,t}(\xi) = \exp\bigl[-\bigl(F(\chi(t)|\xi|) - F(\chi(s)|\xi|)\bigr)\bigr],
  \qquad 0 \le s \le t,\ \xi \in \RR^d .
\end{equation}
In particular the kernels from scale zero are $\widehat\mu_{0,t}(\xi) = e^{-F(\chi(t)|\xi|)}$. The
pair is unique up to $(\chi,F) \mapsto (c\,\chi,\,F(c^{-1}\,\cdot))$ with $c > 0$, in the sense that if a gauge $\chi'$
and a real function $F'$ on $[0,\infty)$ with $F'(0) = 0$ also satisfy
\ref{eq:isotropic-classification}, then
$\chi' = c\,\chi$ and $F' = F(c^{-1}\,\cdot)$ for some $c > 0$, and $c = 1$ when
$\chi(1) = \chi'(1) = 1$.
\end{theorem}
% proof omitted: proved in the isotropic paper, the necessity in its section 3.2, the rest in A.2

The classification reduces to the line through the marginals of a family along a line. Two lemmas
of the isotropic paper carry the reduction, and the proof of Proposition~\ref{prop:isotropic-corner}
uses them too.

% shared with blueprint lem:radial-marginal
\begin{lemma}[a rotation-invariant law is determined by its marginal, {\prog{Iso}{Lem.~3.4}}]\label{lem:radial-marginal}
Let $d \ge 1$, let $e \in \RR^d$ be a unit vector and $\pi(x) := \langle e, x\rangle$, and for a
probability measure $\mu$ on $\RR^d$ write $\pi_*\mu$ for its image on $\RR$.
\begin{enumerate}
\item If $Q_*\mu = \mu$ for every $Q \in \Orth(d)$, then $\kappa := \pi_*\mu$ is a symmetric
probability measure on $\RR$, and
\[
  \widehat\mu(\xi) = \widehat\kappa(|\xi|), \qquad \xi \in \RR^d .
\]
In particular $\widehat\mu$ is real and radial, and $\kappa$ does not depend on $e$.
\item Under the hypothesis of (1), $\pi_*\mu = \delta_0$ if and only if $\mu = \delta_0$.
\item For all probability measures $\mu, \mu'$ on $\RR^d$ and every $\lambda > 0$,
$\pi_*(\mu * \mu') = \pi_*\mu * \pi_*\mu'$, $\pi_*\bigl((\lambda I)_*\mu\bigr) =
(\lambda I)_*(\pi_*\mu)$ on $\RR$, and $\pi_*\delta_0 = \delta_0$.
\end{enumerate}
\end{lemma}
% proof omitted: proved in the isotropic paper, its section A.2

% shared with blueprint lem:isotropic-marginal-family@e81ae6e86684+f75ee0c4dd9e
\begin{lemma}[the marginals of an isotropic family form a family on the line, {\prog{Iso}{Lem.~3.5}}]\label{lem:isotropic-marginal-family}
Let $(\mu_{s,t})$ be an isotropic cascade family in dimension $d$ with dilation action
$(S_\lambda)$, let $e$ and $\pi$ be as in Lemma~\ref{lem:radial-marginal}, and put
\[
  \kappa_{s,t} := \pi_*\mu_{s,t}, \qquad \Phi_{s,t}f := \kappa_{s,t} * f
  \quad (f \in L^1(\RR)), \qquad 0 \le s \le t .
\]
Then $\widehat\mu_{s,t}(\xi) = \widehat\kappa_{s,t}(|\xi|)$ for every $\xi \in \RR^d$, each
$\kappa_{s,t}$ is a symmetric probability measure, and $(\Phi_{s,t})$ is a symmetric cascade
measurement family in the sense of \restated{V}{Def.~3.1}{def:line-family}, with the same $(S_\lambda)$.
\end{lemma}
% proof omitted: proved in the isotropic paper's supplement, its section S3.2

The origin data below describe the laws of the class $\mathcal A_d$ through their \Levy{} measures.
The \Levy{} measure of a rotation-invariant self-decomposable law is described by one function of
the radius, by the polar criterion of self-decomposability (\sscite{sato1999levy}, Thm.~15.10,
p.~95), and \S\ref{sec:affine-stable-cone} reads its statements through that function.

% WR s2 (F3): the label sec:line-bridge below is the one that the shared definition's citation of
% the cone paper's Def. 4.11 reads, so that the definition sends the reader to this subsection,
% where the delay datum is restated.
Some of these laws come from \emph{subordination}, which builds kernels from the Gaussian. Blur
with a Gaussian whose variance is drawn at random, and average over the draws. The result is a
\emph{Gaussian variance mixture}, the law of a Brownian motion at a random time (Bochner's
subordination, \sscite{sato1999levy}, \S30, p.~197). The \Matern{} kernel with transform
$(1 + t^2|\xi|^2/2)^{-\gamma}$ is one: its variance is drawn from the gamma law with Laplace
transform $(1+t^2\sigma)^{-\gamma}$ \prog{Iso}{Prop.~3.12(2)}. The cone paper draws the variance
from the delay laws of the time-causal theory, and describes the law of the variance by a
\emph{delay datum} \prog{B}{Def.~4.11}\label{sec:line-bridge}: a pair $(b_0, k_{\mathrm I})$ with
$b_0 \ge 0$ and $k_{\mathrm I} : (0,\infty) \to [0,\infty)$ measurable, subject to
$\int_0^1 k_{\mathrm I}(u)\,du < \infty$ and $\int_1^\infty k_{\mathrm I}(u)\,u^{-1}du < \infty$,
with exponent
\[
  F_{\mathrm I}(\sigma) = b_0\,\sigma + \int_0^\infty\bigl(1 - e^{-\sigma u}\bigr)\,k_{\mathrm I}(u)\,\frac{du}{u},
  \qquad \sigma \ge 0 .
\]
Its polar profile in dimension $d$ is the integral of \ref{eq:polar-profile-d} below, which is
\prog{B}{(4.6)}, and in Definition~\ref{def:isotropic-origin-datum} the drift $b_0$ is written $a$.
This paper uses the subordinated data as examples of profiles bounded at the origin
(\S\ref{sec:affine-stable-cone}). The isotropic paper names the function of the radius, and the
data that subordination gives, as follows.

% shared with blueprint def:isotropic-origin-datum
\begin{definition}[isotropic origin data in dimension $d$, and the polar profile, {\prog{Iso}{Def.~3.10}}]\label{def:isotropic-origin-datum}
Fix an integer $d \ge 1$ and write $\sigma_{d-1}$ for the surface measure of the unit sphere
$S^{d-1} \subset \RR^d$ and $|S^{d-1}| = \sigma_{d-1}(S^{d-1})$ for its surface area. An
\emph{isotropic origin datum in dimension $d$} is a pair
$(a, k_d)$ with $a \ge 0$ and $k_d : (0,\infty) \to [0,\infty)$ nonincreasing, subject to
\[
  \int_0^1 k_d(r)\,r\,dr < \infty
  \qquad\text{and}\qquad
  \int_1^\infty k_d(r)\,\frac{dr}{r} < \infty .
\]
Its \Levy{} measure is the rotation-invariant measure $\nu_d$ on $\RR^d \setminus \{0\}$ with radial
density $f_d(r) = k_d(r)\,\bigl(|S^{d-1}|\,r^d\bigr)^{-1}$, so that
\begin{equation}\label{eq:polar-decomposition}
  \int_{\RR^d} g \,d\nu_d
   = \frac{1}{|S^{d-1}|}\int_{S^{d-1}}\!\!\sigma_{d-1}(d\eta)\int_0^\infty g(r\eta)\,k_d(r)\,\frac{dr}{r}
\end{equation}
for every nonnegative Borel $g$ vanishing at the origin. Its \emph{exponent} is
\begin{equation}\label{eq:isotropic-exponent}
  F_d(\xi) = \tfrac12 a|\xi|^2
   + \int_{\RR^d}\bigl(1 - \cos\langle \xi, x\rangle\bigr)\,\nu_d(dx), \qquad \xi \in \RR^d ,
\end{equation}
and its \emph{law} is the symmetric infinitely divisible law $\mu$ on $\RR^d$ with
$\widehat\mu = e^{-F_d}$. We call $k_d$ the datum's \emph{polar profile} and
$c := k_d(0{+}) \in [0,\infty]$ its \emph{profile at the origin}. The datum is
\emph{subordinated} when there is a delay datum $(a, k_{\mathrm I})$ in the sense of
\restated{B}{Def.~4.11}{sec:line-bridge}, with $k_{\mathrm I}$ nonincreasing, whose polar profile \prog{B}{(4.6)} is $k_d$:
\begin{equation}\label{eq:polar-profile-d}
  k_d(r) = \frac{2}{\Gamma(d/2)}\int_0^\infty v^{d/2-1}e^{-v}\,
    k_{\mathrm I}\Bigl(\frac{r^2}{2v}\Bigr)\,dv , \qquad r > 0 .
\end{equation}
\end{definition}

An origin datum describes its law, a rotation-invariant symmetric infinitely divisible law, through
the coefficient $a$ of its Gaussian part and the polar profile $k_d$.
% Round 2 (presentation review), item 5 (records/affine/reviews/VERIFY-round2-notation-paper.md § 5):
% the link of the class A_d to the origin data, with a pointer to the released statement, [Iso,
% Prop. 3.11](1)-(2), compared with it (F : [0,inf) -> [0,inf); uniqueness of k_d up to a null set).
The data describe exactly the class $\mathcal A_d$: a function $F : [0,\infty) \to [0,\infty)$ lies in
$\mathcal A_d$ if and only if $\xi \mapsto F(|\xi|)$ is the exponent $F_d$ of an isotropic origin
datum in dimension $d$, and the datum is then unique, $k_d$ up to a Lebesgue-null set
\prog{Iso}{Prop.~3.11}. In this paper it describes the isotropic member
$\mu_I$ of an affine orbit by its profile (\S\ref{sec:affine-stable-cone}). The name refers to the
use the isotropic paper makes of it, at the origin of a kernel.

The later sections use one more statement of the isotropic paper.

% shared with blueprint lem:isotropic-exponent-log
\begin{lemma}[a profile bounded at the origin gives a logarithmic exponent, {\prog{Iso}{Lem.~A.2}}]\label{lem:isotropic-exponent-log}
Let $d \ge 1$, let $(0,k_d)$ be an isotropic origin datum in dimension $d$ with Gaussian
coefficient $0$ and $c = k_d(0{+}) < \infty$, let $F_d$ be its exponent, and put
$C_k := \int_1^\infty k_d(r)\,dr/r$.
\begin{enumerate}
\item \emph{(The exponent grows at most like $c\log|\xi|$.)} For every $\xi$ with $|\xi| \ge 1$,
\[
  F_d(\xi) \le c\log|\xi| + \tfrac{13}{4}\,c + 2C_k ,
  \qquad\text{so}\qquad
  e^{-F_d(\xi)} \ge e^{-13c/4 - 2C_k}\,|\xi|^{-c} .
\]
\item \emph{(The exponent is logarithmic.)} For every $a > 0$,
$F_d(a\xi) - F_d(\xi) \to c\log a$ as $|\xi| \to \infty$. Equivalently, for a unit vector $e$ the
function $g(u) := F_d(\sqrt u\,e)$, $u \ge 0$, satisfies $g(au) - g(u) \to \tfrac c2\log a$ as
$u \to \infty$, for every $a > 0$.
\end{enumerate}
\end{lemma}
% proof omitted: proved in the isotropic paper's supplement, its section S3.3

Clause (2) says $e^{-F_d(\xi)} = |\xi|^{-c}L(|\xi|)$ with $L$ \emph{slowly varying}, that is,
$L(a\rho)/L(\rho) \to 1$ as $\rho \to \infty$ for every $a > 0$. It is the link from a profile
bounded at the origin to an exponent of logarithmic growth, which \S\ref{sec:affine-stable-cone}
uses.

\subsection{The cited facts}\label{sec:cited-facts}

% CHANGED Q-0281 wave 2: blind review W2C-1 (the author's decision of 2026-10-05: Mathlib's route).
% CHANGED Q-0281 wave 2, group F: C4 added, ledger A9 (the author's decision of 2026-10-05), which
% lem:spectral-cones(2) now cites in place of the compound-matrix proof.
% Round 3 (report E, section 3.1; the author's decision D10, PLAN-review-response-paper.md R3.0): C5
% added, ledger A2, the interface node thm:cited-A2 transcribed verbatim; it enters with
% prop:wedge-axis-profiles(2), the whole (GL(d), stable) cell.
The statements of this paper rest on five theorems of the literature that it cites without proving
them, four of probability theory and one of matrix analysis. Each is stated below as it is used,
with the page of its source, and named C1 to C5, C for cited, to keep them apart from the axioms
(A1)--(A8); from here on a cited fact is referred to by its name alone. Inside a cited fact the
letters are its source's: $A$ is a matrix, $\nu$ a \Levy{} measure and $\gamma$ a drift vector, and in
C5 $B$ is a Borel set and $\lambda$ a measure on the unit sphere.

A \emph{null array} is a triangular array of laws in which, for every $\varepsilon > 0$, the largest
mass that an entry of row $n$ puts outside the ball of radius $\varepsilon$ tends to $0$ as $n$ grows. Khintchine's theorem says that the row products of such an array can
only converge to an infinitely divisible law. It is what makes the increments of a cascade
infinitely divisible in \S\ref{sec:classical}.

% shared with blueprint thm:cited-A6
\begin{theorem}[C1: Khintchine's theorem, {\sscite{sato1999levy}, Thm.~9.3, p.~47}]\label{thm:cited-A6}
Let $r_n \ge 1$ be integers and $\mu_{n,k}$, $1 \le k \le r_n$, $n = 1, 2, \dots$, probability
measures on $\RR^d$ with $\lim_{n\to\infty}\max_{1\le k\le r_n}\mu_{n,k}(\{|x| > \varepsilon\}) = 0$ for every
$\varepsilon > 0$. If $\mu_{n,1} * \dots * \mu_{n,r_n} * \delta_{-b_n}$ converges weakly to a
probability measure $\mu$ for some $b_n \in \RR^d$, then $\mu$ is infinitely divisible.
\end{theorem}

Theorem~\ref{thm:cited-A6} is stated in the source for the row sums of independent random
variables. The form above restates it through the law of a sum of independent variables, which is
the convolution of their laws (\sscite{sato1999levy}, p.~8).

The \emph{generating triplet} of an infinitely divisible law is the triple of its
\Levy--Khintchine representation: a Gaussian covariance matrix $A$, a \Levy{} measure $\nu$ and a
drift $\gamma$. The paper compares the kernels of an affine orbit through their \Levy{} measures
(\S\ref{sec:affine-stable-cone}), and C2 is what lets a law and its triplet stand for each other.

% shared with blueprint thm:cited-A3@27418ee4f07a
\begin{theorem}[C2: the \Levy--Khintchine representation, {\sscite{sato1999levy}, Thm.~8.1, pp.~37--38}]\label{thm:cited-A3}
On $\RR^d$:
\begin{enumerate}
\item[(i)] An infinitely divisible law $\mu$ on $\RR^d$ has
\[
  \widehat\mu(z) = \exp\Bigl[-\tfrac12\langle z, Az\rangle + i\langle\gamma, z\rangle
   + \int_{\RR^d}\bigl(e^{i\langle z,x\rangle} - 1 - i\langle z,x\rangle\mathbf{1}_{\{|x| \le 1\}}\bigr)\,
     \nu(dx)\Bigr]
\]
for a symmetric nonnegative definite $A$, a $\gamma \in \RR^d$ and a measure $\nu$ with
$\nu(\{0\}) = 0$ and $\int(|x|^2 \wedge 1)\,\nu(dx) < \infty$.
\item[(ii)] The triplet $(A, \nu, \gamma)$ is unique.
\item[(iii)] Conversely, every such triplet is the generating triplet of an infinitely divisible
law.
\end{enumerate}
\end{theorem}

A law is \emph{operator-self\-de\-com\-pos\-able} when the contraction in the definition of
self\-de\-com\-pos\-abil\-ity is replaced by the linear flow $e^{-tB}$ of a matrix $B$: the law is its own
image under $e^{-tB}$ convolved with a residual, for every $t \ge 0$. The class is Urbanik's
\sscite{urbanik1972levy}. The fact used of it is that, when the flow contracts, $e^{-tB} \to 0$ as
$t \to \infty$, the residuals and the law are infinitely divisible. Without that hypothesis the fact
fails: for $B = 0$ every law is its own image with the residual $\delta_0$.

% shared with blueprint thm:cited-A5
\begin{theorem}[C3: operator-self-decomposable laws with a contracting exponent are infinitely divisible, {\sscite{jurek1993operator}, Thm.~3.3.5, p.~83, with \S3.4, pp.~89--90; \sscite{urbanik1972levy}, Prop.~5.2, p.~134}]\label{thm:cited-A5}
Let $B$ be a real $d \times d$ matrix with $e^{-tB} \to 0$ as $t \to \infty$, and let $\mu$ be a
probability measure on $\RR^d$ such that for every $t \ge 0$ there is a probability measure
$\rho_t$ with $\mu = (e^{-tB})_*\mu * \rho_t$. Then every $\rho_t$ is infinitely divisible, and so
is $\mu$.
\end{theorem}

Jurek and Mason state their Theorem~3.3.5 for \emph{full} laws, whose support lies in no proper
hyperplane (\sscite{jurek1993operator}, p.~46). The opening paragraph of their \S3.4 notes that fullness is used only in the other half of
that theorem and drops it, and adds that the law itself is then infinitely divisible. C3 is the last
sentence of the theorem read with that paragraph. Urbanik's Proposition~5.2 states the same without
fullness. Beyond the letter of the two passages, C3 reads two things. A measure is a probability
measure, which is the book's convention. And the passages say that each residual of one family is
infinitely divisible, where C3 says it of every probability measure that is a residual at some
$t \ge 0$; this is their sentence applied to a family completed around that measure, the hypothesis
supplying the residuals at the other $t$.

The fourth fact is of matrix analysis. The parameter of an aperture in this paper is a positive
semidefinite matrix, and an order on these matrices invariant under the general linear group is determined by the relative steps
$\Sigma^{-1/2}\Sigma'\Sigma^{-1/2}$ that it allows (Proposition~\ref{prop:invariant-orders}). The
composite of two steps depends on how their eigenframes sit relative to each other, and C4 bounds
its spectrum by the spectra of the two (Lemma~\ref{lem:spectral-cones}(2),
Appendix~\ref{app:wedge-axis-orders}).
The eigenvalues of such a product $MN$ are real and positive, being those of $M^{1/2}NM^{1/2}$; write
$\lambda^\downarrow(A)$ for the eigenvalues of a matrix $A$ with real eigenvalues, in decreasing order and
with multiplicity.

% Ledger entry A9, the weak form, re-read at the page on 2026-10-06
% (records/affine/READING-bhatia1997matrix-lidskii.md).
% shared with blueprint thm:cited-A9@ee60aba37e90
\begin{theorem}[C4: the eigenvalues of a product of positive definite matrices, {\sscite{bhatia1997matrix}, Cor.~III.4.6, (III.21), p.~73}]\label{thm:cited-A9}
Let $M$ and $N$ be positive definite $d \times d$ matrices. For $1 \le k \le d$,
\[
  \sum_{i \le k}\log\lambda^\downarrow_i(MN) \le \sum_{i \le k}\bigl(\log\lambda^\downarrow_i(M) +
  \log\lambda^\downarrow_i(N)\bigr) .
\]
\end{theorem}

The source states (III.21) for complex positive matrices as a majorization,
$\log\lambda(AB) \prec \log\lambda^\downarrow(A) + \log\lambda^\downarrow(B)$: the inequalities of C4
between the partial sums of the $k$ largest entries, with the equality of the totals, which is
$\det(MN) = \det M\det N$. The same inequality is stated as Exercise~6.5.3, (6.47), p.~228 of
\sscite{bhatia2007positive}, with a pointer to the page above, as an inequality due to Gel'fand,
Naimark and Lidskii. The notes to Chapter~III of \sscite{bhatia1997matrix} (p.~80) give their papers
of 1950.
% Bhatia's attribution, read at the page images by the librarian on 2026-10-14
% (records/affine/READING-bhatia1997matrix-lidskii.md, section 4); the primary sources are not held.

The fifth fact describes the self-decomposable laws among the infinitely divisible ones by their
\Levy{} measures. In polar coordinates the \Levy{} measure of such a law has, in each direction, a
density with respect to $dr/r$ that does not increase with the radius $r$.

% Ledger entry A2, read at the page (blueprint/AXIOMS.md A2; records/affine/READING-sato1999levy-dimension-d.md).
% shared with blueprint thm:cited-A2
\begin{theorem}[C5: the polar criterion, {\sscite{sato1999levy}, Thm.~15.10, p.~95}]\label{thm:cited-A2}
An infinitely divisible $\mu$ on $\RR^d$ with generating triplet $(A,\nu,\gamma)$ is
self-decomposable if and only if
$\nu(B)=\int_S\lambda(d\xi)\int_0^\infty\mathbf 1_B(r\xi)\,k_\xi(r)\,dr/r$ for a finite measure
$\lambda$ on the unit sphere $S=S^{d-1}$ and a nonnegative $k_\xi(r)$, measurable in $\xi$ and
nonincreasing in $r>0$.
\end{theorem}

Theorem~\ref{thm:cited-A2} is the source's formula (15.8) with both directions, and the source adds that self-decomposability puts no
restriction on $A$ and $\gamma$; its Corollary~15.11, on the same page, is the case $d = 1$.
Self-decomposability is \ref{eq:self-decomposable}, which is the source's Definition~15.1, and the
triplet is that of C2. For a rotation-invariant law the isotropic paper reads C5 as the polar form of
its class (\prog{Iso}{Prop.~3.11}, \S\ref{sec:iso-recap}).

Each enters with a few statements. C1 makes the increments of a cascade infinitely divisible
(Theorem~\ref{thm:cascade-infinitely-divisible}). C2 enters with the restated classification
(Theorem~\ref{thm:isotropic-classification}) and with the laws of the affine orbits and their
\Levy{} measures, from Lemma~\ref{lem:elliptical-stable} on. C3 enters with the anisotropy budget
(Proposition~\ref{prop:anisotropy-budget}(1)) and through it with the theorem on the general linear group along the chain of cones
(Theorem~\ref{thm:wedge-axis}). C4 enters with the spectrum of a product
(Lemma~\ref{lem:spectral-cones}(2)) and through it with the invariant orders
(Theorem~\ref{thm:cone-orders}). C5 enters once, with the description of the whole general linear cell
along a cone order (Proposition~\ref{prop:wedge-axis-profiles}(2)). Appendix~\ref{app:verification} lists every statement that rests on each.

% Q-0309 (fix group 4, the whole-paper review's change 3 and its clash map): a notation table for
% sections 3-4, paper-only. It lists each symbol once, where it is defined; the letters that keep a
% second meaning inside a shared statement are named after it (their renames are owed at the
% blueprint, records/affine/GATE-affine.md).
% WR s2 (S3): the subsection heading is folded into section 2.1, which announces the table.

\begin{table}[ht]
\centering
\small
\setlength{\tabcolsep}{4pt}
\renewcommand{\arraystretch}{1.1}
\begin{tabular}{@{}p{0.27\textwidth}p{0.45\textwidth}p{0.24\textwidth}@{}}
\toprule
symbol & meaning & defined in \\
\midrule
$\Sigma, \Sigma' \in \PSD(d)$; $\PD(d)$ & aperture parameters; the positive definite ones &
Def.~\ref{def:gw-cascade-family} \\
$(G, W)$, $\le_W$, $P_W$ & wedge datum: covariance group, order, its pairs & Def.~\ref{def:gw-cascade-family} \\
$A \cdot \Sigma$, $A_*\mu$ & $A\Sigma A^{\tp}$, and the image of a law & Def.~\ref{def:gw-cascade-family} \\
$\mu_{\Sigma,\Sigma'}$, $\mu_\Sigma$ & kernel of a pair; kernel from scale zero $\mu_{0,\Sigma}$ &
Def.~\ref{def:gw-cascade-family} \\
$\mathrm S(d)$, $\mathrm D_+(d)$, $\GL^+(d)$ & similarities, positive diagonal, orientation-preserving
group & \S\ref{sec:scale-shape} \\
% Round 2 (presentation review), item 20: rows for D_s, Delta(Sigma), sym B and N(0, C), the symbols
% of sections 3-4 the table lacked (VERIFY-round2-structure-paper.md § 20); D_s and Delta(Sigma) are
% defined in Table 2's legend (item 4), sym B in Prop. 4.10, N(0, C) in section 4.9 before its first
% use in Prop. 4.24.
$D_s$; $\Delta(\Sigma)$ & $\diag(s_1,\dots,s_d)$; the diagonal of $\Sigma$ & \S\ref{sec:landscape-table} \\
$g$ & $\widehat\mu_\Sigma(\omega) = e^{-g(\langle\omega,\Sigma\omega\rangle)}$ on an affine orbit &
Lem.~\ref{lem:affine-reduction}, Prop.~\ref{prop:affine-orbits} \\
$M$; $\lambda_{\min}(M)$ & relative step $\Sigma^{-1/2}\Sigma'\Sigma^{-1/2}$; its least
eigenvalue & Prop.~\ref{prop:invariant-orders} \\
$\mathcal K$, $\le_{\mathcal K}$ & step set in $\PD(d)$, and its relation &
Prop.~\ref{prop:invariant-orders} \\
$C$, $\preceq_C$, $\preceq^+_C$ & invariant cone in $\Sym(d)$; its order on $\PD(d)$, on $\PSD(d)$ &
Thm.~\ref{thm:cone-orders} \\
$C_\alpha$, $0 \le \alpha \le \infty$ & the chain of cones from the rays to $\PSD(d)$ &
Prop.~\ref{prop:stable-order-generated} \\
$\preceq_\alpha$; $\le_\alpha = \preceq^+_{C_\alpha}$ & stable order on $\PD(d)$; on $\PSD(d)$ &
Props.~\ref{prop:stable-step-cone}, \ref{prop:similarity-stable-order-families} \\
$\mathcal K_\alpha$ & stable step cone, the step set of $\preceq_\alpha$ &
Thm.~\ref{thm:cascade-infinitely-divisible} \\
$\mathfrak k_\alpha \subseteq \gl(d)$ & exponent cone of $S_\alpha(I)$ & Prop.~\ref{prop:stable-exponent-cone} \\
$\mathrm{sym}\,B$ & $\tfrac12(B + B^{\tp})$, the symmetric part of $B$ & Prop.~\ref{prop:stable-exponent-cone} \\
$S_\alpha(H)$, $n_H$, $\psi_H$ & elliptical stable law, its \Levy{} density and exponent &
Lem.~\ref{lem:elliptical-stable} \\
$k$; $R(\lambda)$; $\beta_-(k)$, $\beta(k)$ & profile, $k = k_d/|S^{d-1}|$; dilation ratio; lower and
growth exponents & Prop.~\ref{prop:orbit-step-cone}, Lem.~\ref{lem:dilation-ratio} \\
$k_d$, $k_d(0{+})$ & polar profile and its value at the origin &
Def.~\ref{def:isotropic-origin-datum} \\
$t(\Sigma)$, $S(\Sigma)$, $N$, $n_i$ & volume scale, shape, relative shape, its eigenvalues &
Lem.~\ref{lem:loewner-scale-shape} \\
$\operatorname{es}$, $\operatorname{ecc}$, $\operatorname{cond}$, $\mathcal M_s$ & exponent spread,
eccentricity, condition number, moment matrix & Props.~\ref{prop:similarity-stable-order-kernels},
\ref{prop:similarity-general-bounded-shape} \\
$N(0, C)$ & the centred Gaussian law of covariance $C$ & \S\ref{sec:smaller-groups} \\
$\mathcal L_\Sigma$ & the operator $-\nabla\cdot\Sigma\nabla$ & \S\ref{sec:affine-stable-cone} \\
\bottomrule
\end{tabular}
\caption{Notation of \S\S\ref{sec:classical}--\ref{sec:landscape}.}
\label{tab:notation}
\end{table}

Some letters keep a second meaning inside a statement, which that statement fixes. They are these.
\begin{itemize}
\item $a$: the Gaussian coefficient $a_{\mathrm V}$ of \ref{eq:admissible-line} on the line, and the
coefficient $a$ of an origin datum, whose exponent has the term $\tfrac12a|\xi|^2$, so that
$a = 2a_{\mathrm V}$ at $d = 1$. ``Gaussian coefficient $0$'' in
Lemma~\ref{lem:isotropic-exponent-log} means $a = 0$, and its clause (2) uses $a > 0$ for a dilation
factor.
\item $c$: the profile at the origin in Definition~\ref{def:isotropic-origin-datum} and
Lemma~\ref{lem:isotropic-exponent-log}, and a positive constant elsewhere.
\item $C$: an invariant cone, with $C(\mathcal K)$ the cone of a step set, and the Gaussian covariance
$C(\Sigma)$ of a kernel in Proposition~\ref{prop:similarity-stable-order-kernels}; a real matrix in
Lemma~\ref{lem:stable-tail}. $C_k$ is a
constant in Lemma~\ref{lem:isotropic-exponent-log}, and C1--C5, upright and without a subscript, are
the cited facts.
\item $g$: the function of one variable through which an affine orbit depends on the aperture,
$\widehat\mu_\Sigma(\omega) = e^{-g(\langle\omega,\Sigma\omega\rangle)}$, in every statement that uses
it. In Lemma~\ref{lem:isotropic-exponent-log}(2) it is the same function for the orbit of a radial
exponent.
\item $\gamma$: a stable or \Matern{} index, and the drift of a triplet in C2.
\item $k$: the folded profile of \ref{eq:admissible-line} on the line, the polar profile $k_d$ of
Definition~\ref{def:isotropic-origin-datum}, and the profile $k$ of
Proposition~\ref{prop:orbit-step-cone}, taken against $r^{-d}dx$, which is $k_d/|S^{d-1}|$. At $d = 1$
the last is half the folded profile, which counts both directions.
\item $S$: the similarities $\mathrm S(d)$, the shape $S(\Sigma)$ of an aperture, the elliptical stable
law $S_\alpha(H)$, the relabelling $S_\lambda$ of the scales in (A8) and (I4), and the unit sphere
$S^{d-1}$, kept apart by their typeface and arguments.
\item $\sigma$: the surface measure $\sigma_{d-1}$, the normalized spectrum $\sigma(\Sigma)$ of
Proposition~\ref{prop:similarity-ray-families}, and, in the prose of \S\ref{sec:shape-adapted}, the
slant of a surface.
\item $t$: the volume scale $t(\Sigma)$ of Lemma~\ref{lem:loewner-scale-shape}, the time of a linear flow
in Proposition~\ref{prop:one-parameter-families}, and the threshold $t_\alpha(\Sigma)$ of
Proposition~\ref{prop:similarity-stable-order-families}(1). The prose writes $v$ and $w$ for variances
and $e^{sB}$ for a flow.
\end{itemize}
Two objects have several names. The stable order of index $\alpha$ is $\preceq_\alpha$ on $\PD(d)$,
defined by its step condition in Proposition~\ref{prop:stable-step-cone}(2).
Proposition~\ref{prop:stable-order-generated}(1) shows that it equals $\preceq_{C_\alpha}$. On
$\PSD(d)$ it is $\le_\alpha$, which equals $\preceq^+_{C_\alpha}$. Its set of steps is the step cone
$\mathcal K_\alpha = \exp C_\alpha \subseteq \PD(d)$, a set of relative steps $M$. The exponent cone
$\mathfrak k_\alpha \subseteq \gl(d)$ is a set of exponents, in fraktur as $\gl(d)$ is, and its
symmetric part is $C_\alpha$ (Proposition~\ref{prop:stable-order-generated}(2)).